\documentclass[tikz,border=6pt]{standalone}
\usepackage{ctex}
\usepackage{amsmath}
\usetikzlibrary{calc, arrows.meta}

\begin{document}
\begin{tikzpicture}[
  scale=1.1, thick,
]

% ── 垂直特例：l1 水平（k1=0），l2 竖直（k2 不存在）──
% 坐标轴（浅灰色参考）
\draw[gray!40, thin, ->] (-0.3, 0) -- (4.2, 0) node[right, font=\footnotesize, gray] {$x$};
\draw[gray!40, thin, ->] (0, -0.3) -- (0, 3.5) node[above, font=\footnotesize, gray] {$y$};

% l1：水平线 y = 1.5
\draw[black, thick] (-0.1, 1.5) -- (3.9, 1.5);
\node[font=\small] at (4.05, 1.5) {$l_1$};
\node[font=\footnotesize, black] at (1.8, 1.75) {$k_1 = 0$（水平）};

% l2：垂直线 x = 2.2
\draw[red, thick] (2.2, -0.1) -- (2.2, 3.3);
\node[font=\small, red] at (2.2, 3.5) {$l_2$};
\node[font=\footnotesize, red] at (3.2, 2.8) {$k_2$ 不存在};

% 直角标记（交点 (2.2, 1.5)）
\draw[gray, thin]
  (2.2, 1.5) --
  (2.2 + 0.18, 1.5) --
  (2.2 + 0.18, 1.5 + 0.18) --
  (2.2, 1.5 + 0.18);

% 交点圆点
\filldraw[black] (2.2, 1.5) circle (1.5pt);

% 警告文字框
\draw[red!60, rounded corners=4pt, thick]
  (0.1, -0.9) rectangle (3.9, -0.1);
\node[font=\footnotesize, red, align=center] at (2.0, -0.5)
  {$k_1 k_2 = -1$ \textbf{不能直接用}！\\
   应用一般式：$A_1 A_2 + B_1 B_2 = 0$};

\end{tikzpicture}
\end{document}
